% !TEX root = ./main.tex

% 工作标题及个人信息。正式提交前核对学院规定并填写真实信息。
\sjtusetup{
  info = {%
    zh / title      = {面向关系表格学习的语义建模与高效训练方法研究},
    en / title      = {Semantic Modeling and Efficient Training Methods for Relational Table Learning},
    zh / keywords   = {关系表格学习, 元路径, 异构图, 按需训练},
    en / keywords   = {relational table learning, meta-path, heterogeneous graph, on-demand training},
    zh / author     = {待填写},
    en / author     = {To be completed},
    id              = {待填写},
    zh / supervisor = {待填写},
    en / supervisor = {To be completed},
    zh / department = {待填写},
    en / department = {To be completed},
    zh / major      = {待填写},
    en / major      = {To be completed},
    zh / degree     = {待填写},
    en / degree     = {To be completed},
    zh / display-date = {待填写},
  },
}

% 使用 BibLaTeX 处理参考文献
%   biblatex-gb7714-2015 常用选项
%     gbnamefmt=lowercase     姓名大小写由输入信息确定
%     gbpub=false             禁用出版信息缺失处理
\usepackage[backend=biber,style=gb7714-2015]{biblatex}
% 文献表字体
\renewcommand{\bibfont}{\zihao{5}\setbaselineskip{16bp}}
% 文献表条目间的间距
\setlength{\bibitemsep}{3bp plus 1pt}
% 导入参考文献数据库
\addbibresource{refs.bib}

% 脚注格式
\usepackage[perpage,bottom,hang]{footmisc}

% 定义图片文件目录与扩展名
\graphicspath{{figures/}}
\DeclareGraphicsExtensions{.pdf,.eps,.png,.jpg,.jpeg}

% 颜色宏包
\usepackage{xcolor} 

% 写作阶段的可见占位文字；定稿时逐段替换。
\newcommand{\thesisplaceholder}[1]{%
  \par\smallskip\noindent{\color{gray}\itshape 待补充：#1}\par\smallskip
}

% 确定浮动对象的位置，可以使用 [H]，强制将浮动对象放到这里（可能效果很差）
% \usepackage{float}

% 固定宽度的表格
% \usepackage{tabularx}

% 使用三线表：toprule，midrule，bottomrule。
\usepackage{booktabs}

% 表格中支持跨行
\usepackage{multirow}

% 表格中数字按小数点对齐
\usepackage{dcolumn}
\newcolumntype{d}[1]{D{.}{.}{#1}}

% 使用长表格
\usepackage{longtable}

% 附带脚注的表格
\usepackage{threeparttable}

% 附带脚注的长表格
\usepackage{threeparttablex}

% 算法环境宏包
\usepackage[ruled,vlined,linesnumbered]{algorithm2e}
% \usepackage{algorithm, algorithmicx, algpseudocode}

% 代码环境宏包
\usepackage{listings}
\lstdefinestyle{lstStyleCode}{%
  aboveskip         = \medskipamount,
  belowskip         = \medskipamount,
  basicstyle        = \ttfamily\zihao{-5}\setbaselineskip{12bp},
  commentstyle      = \slshape\color{black!60},
  stringstyle       = \color{green!40!black!100},
  keywordstyle      = \bfseries\color{blue!50!black},
  extendedchars     = false,
  upquote           = true,
  tabsize           = 2,
  showstringspaces  = false,
  xleftmargin       = 1em,
  xrightmargin      = 1em,
  breaklines        = false,
  framexleftmargin  = 1em,
  framexrightmargin = 1em,
  backgroundcolor   = \color{gray!10},
  columns           = flexible,
  keepspaces        = true,
  % texcl             = true,  % 代码中的注释部分将会被视为 LaTeX 正文模式内容
  % mathescape        = true,  % 代码中$..$的部分将会被视为 LaTeX 数学模式内容
}
\lstnewenvironment{codeblock}[1][]{%
  \lstset{style=lstStyleCode,#1}}{}

% 直立体数学符号
\providecommand{\dd}{\mathop{}\!\mathrm{d}}
\providecommand{\ee}{\mathrm{e}}
\providecommand{\ii}{\mathrm{i}}
\providecommand{\jj}{\mathrm{j}}

% 国际单位制宏包
\usepackage{siunitx}

% 定理环境宏包
\usepackage{amsthm}
% \usepackage{ntheorem}

% 绘图宏包
\usepackage{tikz}
\usetikzlibrary{arrows.meta, shapes.geometric}

% 数据图表宏包
\usepackage{pgfplots}
\pgfplotsset{compat=newest}

% 一些文档中用到的 logo
\usepackage{hologo}
\providecommand{\XeTeX}{\hologo{XeTeX}}
\providecommand{\BibLaTeX}{\textsc{Bib}\LaTeX}

% 借用 ltxdoc 里面的几个命令方便写文档
\DeclareRobustCommand\cs[1]{\texttt{\char`\\#1}}
\providecommand\pkg[1]{{\sffamily#1}}

% hyperref 宏包在最后调用
\usepackage{hyperref}

% E-mail
\providecommand{\email}[1]{\href{mailto:#1}{\urlstyle{tt}\nolinkurl{#1}}}
